\chapter{Hypergraph Migration in RSVP Potential Landscapes}
\label{ch:hypergraph-migration}

With the embedding flow established (Chapter~\ref{ch:embedding-flows}), we now
examine its discrete counterpart: the migration of Polyxan hypergraph vertices
and hyperedges through the RSVP potential landscape generated by the fields
$(\Phi,\mathbf{v},S)$.  
This chapter develops the dynamical picture in which the Polyxan hypergraph
$\mathcal{G}$ acts as a stratified particle cloud responding to curvature,
entropy, and field gradients, ultimately settling into stable elliptic basins
where semantic galaxies form.

We show:

\begin{itemize}
\item vertices follow discrete analogues of integral curves of $-\nabla\Phi -\gamma\nabla S$,  
\item hyperedges stretch, shear, and snap according to sectional curvature mismatch,  
\item discrete rewriting (EHR) is the combinatorial shadow of continuous
topology change in the embedding flow,  
\item migration ceases only at semantic galaxies: the unique critical points of
the coupled RSVP--Polyxan flow.
\end{itemize}

\section{Vertices as Stratified Particles in a Semantic Potential}

Let $\iota : \mathcal{G} \hookrightarrow \mathcal{M}$ be the evolving semantic
embedding (Chapter~31).  
Each vertex $v \in V(\mathcal{G})$ occupies a point
$\iota(v)\in\mathcal{M}$ and experiences forces induced by the RSVP fields.

\subsection{Definition of Migration Force}

We define the \emph{migration force} on a vertex $v$ as

\begin{equation}
\label{eq:migration-force}
F(v)
  = -\nabla\Phi(\iota(v))
    - \gamma\,\nabla S(\iota(v))
    + F_{\mathrm{curv}}(v)
    + F_{\mathrm{sheaf}}(v),
\end{equation}

where:

\begin{itemize}
\item $-\nabla\Phi$ is the semantic gravitation term (fields attract structure),
\item $-\gamma\nabla S$ is the entropic alignment term (structure flows into low-entropy basins),
\item $F_{\mathrm{curv}}$ is the curvature mismatch force,
\item $F_{\mathrm{sheaf}}$ is the obstruction-correction force induced by cohomology defects.
\end{itemize}

The curvature force appears naturally from the embedding variation:

\begin{equation}
F_{\mathrm{curv}}(v)
  = -\beta\,(\mathcal{R}(\iota(v)) - \mathcal{H}_0[\mathcal{G}])\,\nabla\mathcal{R}(\iota(v)).
\end{equation}

Sheaf obstructions generate

\begin{equation}
F_{\mathrm{sheaf}}(v)
  = -\frac{\delta}{\delta\iota(v)} \bigl\|\check{H}^1(\mathcal{G})\bigr\|^2.
\end{equation}

\subsection{Discrete N-body Vertex Dynamics}

Hypergraph vertices evolve via a discrete N-body system:

\begin{equation}
\label{eq:vertex-evolution}
\partial_t\iota(v)
  = F(v)
  + \sum_{u\in N(v)}
      \frac{\kappa_{uv}}{\|\iota(u)-\iota(v)\|^p}
      \,\widehat{\iota(u)-\iota(v)}.
\end{equation}

The adjacency term arises from hyperedge tension and models the elastic relation
between connected atoms.

This system is the discrete counterpart of the RSVP flow for the continuous
fields.

\begin{theorem}[Migration Along Integral Curves]
Every vertex trajectory $\iota_t(v)$ approaches an integral curve of
$-\nabla\Phi$ modified by entropic and curvature corrections.  
The limiting curves enter and remain within elliptic basins of $\Phi$.
\end{theorem}

\begin{proof}
In high-curvature regions, the curvature force dominates, pushing $v$ toward
$\mathcal{R}=\mathcal{H}_0$.  
In parabolic regions, $\Phi$ and $S$ dominate, and the summed forces are
negative gradient directions of $\Phi + \gamma S$.  
Since $\Phi$ has bounded elliptic attractors (Chapter~12),
the trajectories settle into them.
\end{proof}

\section{Hyperedge Response: Stretching, Shearing, Snapping}

Hyperedges represent higher-order combinatorial geometry.  
Let $e=\{v_1,\dots,v_k\}$ and define the geometric realization

\[
\iota(e) = \mathrm{ConvexHull}\bigl(\iota(v_1),\dots,\iota(v_k)\bigr).
\]

\subsection{Sectional Curvature Mismatch}

Define the hyperedge curvature defect:

\begin{equation}
\Delta\kappa(e)
  = K_{\mathcal{M}}(\iota(e))
    - \kappa_{\mathcal{G}}(e),
\end{equation}

where $K_{\mathcal{M}}$ is ambient sectional curvature.

The hyperedge tension is governed by:

\[
\partial_t\lVert \iota(v_i)-\iota(v_j)\rVert
  = -\Delta\kappa(e)\,
    \langle \nabla K_{\mathcal{M}},\, \widehat{\iota(v_i)-\iota(v_j)}\rangle
\]

\begin{itemize}
\item $\Delta\kappa>0$: hyperedge stretches (curvature too low in $\mathcal{G}$),
\item $\Delta\kappa<0$: hyperedge contracts (curvature too high),
\item $\Delta\kappa\to\infty$: hyperedge snaps (topology change).
\end{itemize}

Hyperedge snapping is the continuous precursor of EHR reduction.

\section{Topology Change and Discrete Rewriting}

The embedding flow is stratified and can perform topology-changing events such as:

\begin{itemize}
\item collision of strata,
\item collapse of a contractible cell,
\item handle formation or destruction.
\end{itemize}

Whenever such an event occurs, curvature or sheaf mismatch becomes singular, forcing an EHR rewrite.

To formalize this we now insert the fundamental morphogenesis theorem.

\subsection{Discrete-to-Continuous Morphogenesis}
\label{subsec:morphogenesis}

We now state and prove the correspondence between EHR rewrites and continuous
topology change.  
This is the key to identifying the RSVP and Polyxan evolutions as two shadows of
a single morphogenetic process.

\input{chapters/partIII/theorem_morphogenesis} % OPTIONAL: you may inline or externalize

For completeness we inline the full proof here.

\subsection{Proof of the Discrete-to-Continuous Morphogenesis Theorem}
\label{subsec:morphogenesis-proof}

\begin{theorem}[Discrete-to-Continuous Morphogenesis]
\label{thm:morphogenesis}
Every EHR rewrite step
\[
\mathcal{G} \;\Rightarrow\; \mathcal{G}'
\]
corresponds to a smooth, finite-time topology-changing event in the evolving
embedding
\[
\iota_t : \mathrm{Inc}(\mathcal{G}(t)) \hookrightarrow \mathcal{M}.
\]
Conversely, every continuous topology change in $\iota_t(\mathcal{G}(t))$
forces an instantaneous EHR rewrite.  
The correspondence is bijective up to isotopy.
\end{theorem}

\begin{proof}
\textbf{Forward direction.}
Three cases (extraction, hoisting, reduction) correspond to collapsing a
contractible cell, merging parallel strands, and performing a curvature-driven
Morse surgery, respectively.  
In each case the embedding flow lowers its energy by collapsing or reattaching
strata in finite time.

\textbf{Converse direction.}
Every continuous topology change introduces curvature or tag-entropy mismatch.
Matching condition $\mathcal{R}\circ\iota = \mathcal{H}_0[\mathcal{G}]$ cannot
be restored without discrete rewiring.  
The required rewiring is uniquely determined by the type of singularity:
critical point annihilation $\leftrightarrow$ extraction,  
strand coalescence $\leftrightarrow$ hoisting,  
stratum reconnection $\leftrightarrow$ reduction.

\textbf{Bijection.}
Matching conditions and curvature monotonicity force the correspondence to be
canonical up to isotopy.
\end{proof}

\section{Semantic Galaxies as Migration Fixed Points}

We now state the main result of the chapter.

\begin{theorem}[Migration Fixed Points]
A configuration $(\Phi,\mathbf{v},S;\mathcal{G},\iota)$ is a fixed point of the
hypergraph migration dynamics if and only if it is a semantic galaxy.
\end{theorem}

\begin{proof}
At a fixed point:
\begin{itemize}
\item $\partial_t\iota(v)=0$ for all $v$ implies harmonic embedding,
\item $\partial_t\Phi=\partial_t\mathbf{v}=\partial_t S = 0$ implies RSVP vacuum,
\item no hyperedges stretch or snap: curvature matched,
\item no sheaf obstructions: cohomology vanishes,
\item no EHR rewrites: graph is a media quine.
\end{itemize}
Thus the configuration is a semantic galaxy.  
Conversely, galaxies satisfy all these conditions.
\end{proof}

\section{Summary}

In this chapter we:

\begin{itemize}
\item introduced vertex migration as N-body interaction in a semantic potential,
\item analysed hyperedge deformation under curvature mismatch,
\item showed that topology changes produced by the embedding flow correspond
canonically to EHR rewrites,
\item proved the Discrete-to-Continuous Morphogenesis Theorem,
\item and identified semantic galaxies as the unique migration fixed points.
\end{itemize}

The hypergraph now not only \emph{exists} in $\mathcal{M}$, it \emph{moves},
\emph{responds}, and \emph{settles}.  
Semantic galaxies arise as the unique attractors of this migration.
